\documentclass[11pt,letterpaper]{article}
\usepackage[T1]{fontenc}
\usepackage{lmodern}
\usepackage[margin=0.92in]{geometry}
\usepackage{amsmath,amssymb,amsthm,mathtools}
\usepackage{booktabs,longtable,tabularx,array,multicol}
\usepackage{microtype,xcolor,enumitem,xurl}
\usepackage[colorlinks=true,linkcolor=blue!45!black,citecolor=blue!45!black,
  urlcolor=blue!45!black]{hyperref}
\usepackage{fancyhdr}
\pagestyle{fancy}\fancyhf{}
\fancyhead[L]{\small U208: a reduction using entry context}
\fancyhead[R]{\small 29 October 2026}
\fancyfoot[C]{\thepage}
\setlength{\headheight}{14pt}
\setlength{\parindent}{0pt}
\setlength{\parskip}{5pt plus 1pt}
\setlist{itemsep=2pt,topsep=4pt}
\setcounter{tocdepth}{1}
\newtheorem{theorem}{Theorem}[section]
\newtheorem{lemma}[theorem]{Lemma}
\newtheorem{corollary}[theorem]{Corollary}
\newtheorem{proposition}[theorem]{Proposition}
\newcommand{\PA}{\mathrm{PA}}
\newcommand{\Con}{\operatorname{Con}}
\newcommand{\Halt}{\operatorname{Halt}}
\newcommand{\Sub}{\operatorname{Sub}}
\newcommand{\IR}{\mathrm{IR}}
\newcommand{\act}[3]{(#1,\mathrm{#2},#3)}
\newcommand{\art}[1]{\path{#1}}
\newcommand{\old}{M_{209}}
\newcommand{\newm}{M_{208}}
\hypersetup{pdftitle={U208: eliminating a state using entry context},
  pdfauthor={Research prepared for Vladimir Reshetnikov},
  pdfsubject={Comparison of six U209 proposals, an explicit 208-state machine,
  and a proof of preservation of halting}}

\title{\vspace{-0.7cm}\textbf{U208: eliminating a state\\using entry context}\\[8pt]
\large Comparing six proposed machines for \(\Con(\PA)\)\\
and constructing an explicit 208-state candidate}
\author{Research prepared for Vladimir Reshetnikov}
\date{29 October 2026}
\begin{document}
\maketitle\thispagestyle{empty}

\begin{abstract}
The six U209 archives in the supplied ProveIt snapshot describe only two
machines up to rooted, symbol- and direction-preserving state renaming.
Four executable packages and the PDF-only package belong to one class;
the Compatibility package belongs to the other. All executable versions
preserve the same 3,909-primitive program and 5,083-arrow control graph.

We construct a deterministic, binary, one-tape machine with
\textbf{208 working states and 416 transitions}, including blank initialization.
Starting with the Joint U209 table, delete native state 172, redirect its
incoming edge to 211, and replace the zero transition of 211 by a left move
to 217. The justification is a conditional invariant missed by the earlier
single-predecessor-tag searches: every reachable old 211 configuration
reading zero has a one immediately to its left. Every old 172 entry has
the two preceding bits 10. Four exact local identities then yield a
positive block simulation, with one, three, or seven new steps per old step.
The entire tape and head position agree at synchronization points.

The new theorem proves halting equivalence from state 0 on every common
initial tape. It is directly formalizable in PA for numerically coded
finite computations. Composing with the retained U209 logical calibration
gives the conventional mathematical endpoint
\(\PA\vdash\neg\Halt(\newm,\mathrm{blank})\leftrightarrow\Con(\PA)\).
Fresh symbolic replay verifies all 5,084 primitive contracts and all 416
literal rows. The complete startup trace is unchanged. The archive contains
the exact tables, a deterministic builder, an independent local checker,
comparison witnesses, reproducible discovery, and the retained parent.
No new end-to-end proof-assistant certificate is claimed.
\end{abstract}

\begin{center}
\fbox{\begin{minipage}{0.93\linewidth}
\textbf{Principal executable:} \art{machine/u208.tm}. Working states
\(0,\ldots,207\); start 0; separate halt label 208; alphabet \(\{0,1\}\);
all-zero two-way infinite tape; head at coordinate 0; only L/R moves.\par
\textbf{Dense table SHA-256}\par
\small\path{a2fee9e41311c0d43aaac871ea411be24eada17e6b3633c55f40b437101144e1}
\end{minipage}}
\end{center}
\clearpage\tableofcontents\clearpage

\section{Result, scope, and evidence}
\label{sec:result}

The requested threshold, fewer than 209 working states, is achieved by the
principal 208-state candidate. Its complete dense table is printed in
Appendix~\ref{app:table}. A second 208-state table, obtained from the distinct
Compatibility parent, is supplied as a supplementary construction.

Three claims should be distinguished.
\begin{enumerate}
\item \textbf{New preservation theorem.} The principal U208 table and the
exact Joint U209 table have the same halting behavior from state 0 on every
common binary tape. Section~\ref{sec:proof} gives the complete proof from
their literal transitions. It does not depend on the theory U or on the
meaning of the source program.
\item \textbf{Direct implementation check.} Fresh symbolic induction
checks the principal U208 table against every primitive of the unchanged
source, for arbitrary natural register values and arbitrary allowed
spare-register lists. Exact startup replay supplies the blank-input bridge.
These checks corroborate the new preservation theorem and its implementation.
\item \textbf{Consistency endpoint.} The supplied U209 proof-search and
interpretation arguments identify blank halting with a finite PA refutation.
The new preservation theorem transfers that endpoint. Those substantial
inherited mathematical arguments are retained and audited, rather than
silently replaced by a K215 theorem about a different theory and table.
\end{enumerate}

\begin{theorem}[Relative arithmetic endpoint]
\label{thm:relative}
For the exact principal tables delivered here,
\[
\PA\vdash\Halt(\newm,\mathrm{blank})
\leftrightarrow\Halt(\old,\mathrm{blank}).
\]
The proof uses fixed predecessor inventories, four constant-length literal
traces, and induction on finite computations.
\end{theorem}

\begin{corollary}[Endpoint after the retained calibration]
\label{cor:endpoint}
Composing Theorem~\ref{thm:relative} with the mathematical calibration in
the supplied Joint U209 ancestry gives
\[
\PA\vdash
\neg\Halt(\newm,\mathrm{blank})
\leftrightarrow\Con(U)
\leftrightarrow\Con(\PA).
\]
\end{corollary}

This packet supplies an ordinary mathematical proof and executable
certificates. It supplies no newly compiled Lean, Rocq, MM0, or first-order
PA endpoint proof. It does not assert \(\Con(\PA)\), unconditional
nonhalting, or global minimality. The arithmetic equivalence is the statement
to be proved; a consistency assumption is not used to prove that equivalence.

\section{Comparison of the six incoming proposals}
\label{sec:comparison}

\subsection{A fixed source snapshot}

All six archives and the unary-theory archive~\cite{incoming} were retrieved from
\art{VladimirReshetnikov/ProveIt}, directory \art{docs/incoming}, at commit
\begin{center}\small
\path{343da82d7384d10ec9fd1c17361fc923b2d1d117}.
\end{center}
The source listing, exact URLs, archive sizes, Git blob identities, and
SHA-256 hashes are recorded in \art{baseline/source_manifest.json}.
Downloaded bytes were checked against the listed Git blob hashes. The
comparison uses actual transition rows, not titles or declared state counts.

\begin{table}[ht]
\centering\small
\begin{tabularx}{\linewidth}{@{}lXcl@{}}
\toprule
Short name & Archive stem following \texttt{U209\_ConPA} & Class & Dense hash prefix\\
\midrule
Plain & \path{_2026-10-29} & A & \texttt{65ccb606c8e970ad}\\
Compatibility & \path{_Compatibility_2026-10-29} & B & \texttt{bda91bd40880f96e}\\
Compatible Folds & \path{_Compatible_Folds_2026-10-29} & A & \texttt{5ded628d692ff442}\\
Composed & \path{_Composed_2026-10-29} & A & \texttt{3a69998d818b6e63}\\
Joint & \path{_Joint_2026-10-29} & A & \texttt{5ded628d692ff442}\\
Joint (1) & \path{_Joint_2026-10-29 (1)} (PDF only) & A & \texttt{65ccb606c8e970ad}\\
\bottomrule
\end{tabularx}
\caption{All six archives have the extension \texttt{.zip}. There are two
rooted transition-table isomorphism classes. Full hashes are in the receipts.}
\label{tab:comparison}
\end{table}

Each dense table has working states 0 through 208, separate halt 209, and
418 transitions. All 209 working states are reachable in the directed
transition graph from 0. This graph statement does not say that each state
is reached in the particular blank-tape computation.

\subsection{Exact equality and state-renaming witnesses}

A rooted deterministic binary transition table has a canonical breadth-first
enumeration: start with state 0, visit the read-0 edge before the read-1 edge,
and assign a new canonical label on first encounter. Retain write bits,
directions, and the distinguished halt. Equality of the resulting canonical
tables is equivalent to a root- and halt-preserving state bijection because
all working states are graph-reachable. The standalone comparison program
also constructs the forced bijection directly and checks all 418 rows.
For unequal classes it records an explicit conflict in that forced map.
Thus the result is stronger than a comparison of file hashes.

Compatible Folds and Joint have \emph{byte-identical} native and dense
tables. Plain maps to Joint by the native-label cycles
\[
(217\ 129\ 214),\qquad(215\ 204\ 216),
\]
fixing every other label. Here the first cycle means
\(217\mapsto129\mapsto214\mapsto217\). Composed maps to Joint by
\[
(217\ 129\ 204\ 214\ 216),
\]
again fixing all other labels. These are verified transition-by-transition.
No reflection of the tape, exchange of symbols, or altered start state is used.

The sixth archive contains only one PDF. We extracted all 418 transition
rows from its appendix and reconstructed exactly Plain's dense file,
including its hash. Its machine identity therefore does not rely on the
hash printed in the prose. It is an alternative exposition of class A,
not a third executable construction.

\subsection{The common source program}

The five executable packages have the same ordered instruction word:
\[
\#\mathrm{CI}=1558,\quad\#\mathrm{CD}=1175,\quad
\#\mathrm{CB}=80,\quad\#\mathrm{CC}=1095,\quad\#\mathrm{CH}=1.
\]
The total is 3,909 primitive sites. Resolving both outcomes of each
conditional decrement gives 5,083 nonhalt control arrows. Two independently
implemented resolvers, one based on subtree intervals and one on binary
addresses and ancestor cuts, agree on the full ordinal control graph.

All class A source files are identical. They have maximum address depth 42;
Compatibility's different association has maximum depth 41. Both have
197 distinct syntactic subtrees and maximum consecutive-zero run 6. The
depth difference is a layout fact; it is not a proved asymptotic speedup.
All five executable packages retain the same unary-theory ZIP, byte-for-byte. The proposals
therefore concern decoder organization and physical state sharing, not
different calibrating theories.

\subsection{What each presentation contributes}

All five full packages combine the former U211 savings associated with
native states 180, 181, and 218. These savings cannot be added a second time
by combining the current reports.

\textbf{Plain} presents thirty closed-factor rotations from U212, an
arithmetic-entry repair, and three simultaneous deletions. Its useful
additional research is a logical obstruction to deleting the dominated
propositional row; see Section~\ref{sec:limits}.

\textbf{Composed} obtains a renaming of Plain's machine and gives a clear
account of shared helper entries and of the reverse direction of halting
simulation. It also directly retains recovered U222 calibration material.

\textbf{Compatible Folds} uses seventeen rotations, an eleven-factor
reassociation of the emitter, and three outer rotations. Its final machine
is exactly Joint's, despite a different construction history.

\textbf{Joint} gives a compact simultaneous layout of eighteen selectors
over a closed 3,096-primitive region, then performs the three deletions.
Its deterministic builder, direct local checker, immutable engine binding,
nonmutating replay, and complete nested ancestry make it the most convenient
baseline for the present literal-table reduction. That choice is about
auditability and packaging, not a smaller state count than its isomorphic peers.

\textbf{PDF-only Joint (1)} gives twenty-four reassociations from Contextual
U211, corresponding to thirty when the parent's six are included. Its
transition appendix is useful, but the archive itself contains no runnable
builder or verification suite.

\textbf{Compatibility} is the sole different machine. It preserves a
backward call through helper 129 by keeping the rows that call can actually
read, while changing an unused branch, 204/read-1. An eight-factor
reassociation then prepares the fold \(180\mapsto216\). This is a useful
example of preserving the relevant behavior of a shared helper without
freezing its entire table. It is an alternative way to combine the same
three savings, not a fourth independent deletion.

\subsection{Fresh replay of both classes}

One complete release from each class was freshly replayed. Joint's replay
reconstructed its sealed machine and source, checked its local identities,
and rederived all universal primitive contracts. Compatibility's replay
likewise reconstructed both retained stages and rederived its final
contracts, with fresh case records equal to its saved records.

\begin{table}[ht]
\centering
\begin{tabular}{@{}lrr@{}}
\toprule
Check & Joint & Compatibility\\
\midrule
Working states & 209 & 209\\
Universal primitive contracts & 5,084 & 5,084\\
Literal rows covered & 418/418 & 418/418\\
Symbolic inference applications & 1,250,049 & 1,248,748\\
Exact blank startup steps & 48,140 & 48,140\\
Full release replay & PASS & PASS\\
\bottomrule
\end{tabular}
\caption{Fresh verification of the two distinct input designs. These are
executable mathematical inductions in inherited Python engines.}
\end{table}

\section{The exact 208-state construction}
\label{sec:construction}

The parent \(\old\) is Joint's sparse native table, with SHA-256
\begin{center}\small
\path{d0a155d3c683789f3bf5327c0cfb2a0cfd78076af8f1654d4ca426ccbb604d16}.
\end{center}
Native labels are historical names. In particular native 211 is a working
state, while native 254 is halt. They are not counts.

Define the projection
\[
\rho(172)=211,\qquad\rho(q)=q\quad(q\ne172).
\]
Delete both rows of native 172. Redirect every target 172 to 211; there is
exactly one such incoming row. Finally change the read-zero row of 211.
The complete difference in retained rows is
\begin{center}
\begin{tabular}{@{}ccll@{}}
\toprule
State & Read & Old action & New action\\
\midrule
116 & 0 & \(\act{0}{R}{172}\) & \(\act{0}{R}{211}\)\\
211 & 0 & \(\act{0}{R}{157}\) & \(\act{0}{L}{217}\)\\
\bottomrule
\end{tabular}
\end{center}
No other retained transition is edited. In particular,
\(\delta(211,1)=\act{1}{R}{184}\) remains unchanged. The source tree,
primitive order, operands, source addresses, header convention, and register
program are all unchanged.

The remaining native states are sorted and assigned dense labels
0 through 207; native halt 254 becomes dense halt 208.
\art{machine/state_map.json} records the complete bijection.
The deterministic builder checks target closure, binary completeness,
all 208 graph-reachable working states, and both expected file hashes.

\section{Complete proof of the new reduction}
\label{sec:proof}

\subsection{Configurations and the relation}

A configuration is \((q,h,t)\), where \(q\) is a control state,
\(h\in\mathbb Z\), and \(t:\mathbb Z\to\{0,1\}\) is the tape. An action
\(\act{w}{D}{z}\) writes \(w\), moves by \(-1\) or \(+1\), and enters
\(z\). Reaching the distinguished halt ends the computation.
The relation at simulation boundaries is
\[
(q,h,t)\ \mathcal R\ (\rho(q),h,t).
\]
The old state is retained in this relation. In particular old 172 and old
211 have the same new representative but different entry guarantees.

\subsection{Complete predecessor inventories}

The following are the entire incoming inventories in \(\old\), checked
against every one of its 418 rows. The notation \(p/b:wDz\) records
the source state and read symbol as well as the action.
\begin{align*}
\operatorname{In}(172)&=\{116/0:0R172\},\\
\operatorname{In}(116)&=\{215/1:1R116\},\\
\operatorname{In}(211)&=\{163/0:0L211,\ 163/1:1R211\},\\
\operatorname{In}(163)&=\{132/1:1R163,\ 199/1:1R163,\ 217/0:0R163\},\\
\operatorname{In}(217)&=\{216/1:1R217\}.
\end{align*}
None of 172, 116, 211, 163, or 217 is the initial state 0.
Consequently every occurrence needed in the following backward arguments
has the stated predecessor; there are no initial-time exceptions.

\begin{lemma}[Two bits before state 172]
\label{lem:172}
In every computation of \(\old\) beginning in state 0, on any initial
tape, an occurrence of state 172 at head \(h\) satisfies
\(t(h-2)=1\) and \(t(h-1)=0\).
\end{lemma}
\begin{proof}
Its last transition must be 116/read-0, which writes zero at \(h-1\)
and moves right. The preceding transition into 116 must be 215/read-1,
which writes one at \(h-2\) and moves right. The final transition does
not revisit \(h-2\). Both facts therefore hold at entry to 172.
\end{proof}

\begin{lemma}[A conditional bit before state 211]
\label{lem:211}
In every such old computation, an occurrence of state 211 at head \(h\)
with \(t(h)=0\) satisfies \(t(h-1)=1\).
\end{lemma}
\begin{proof}
If the last transition is 163/read-1, it has just written one at \(h-1\)
and moved right, proving the conclusion.

Otherwise the last transition is 163/read-0 from head \(h+1\), moving
left to \(h\). Consider the preceding entry to 163. All three incoming
edges to 163 move right from \(h\). Entries through 132/read-1 or
199/read-1 would have written one at \(h\), and the following move from
\(h+1\) would not change it. Both are incompatible with the present
read-zero assumption. Thus the entry to 163 must have been 217/read-0.
The sole entry to that occurrence of 217 was 216/read-1, which wrote one
at \(h-1\). The next two transitions touch \(h\) and \(h+1\), so
the one at \(h-1\) persists.
\end{proof}

The second lemma is conditional on the current read symbol. The unguarded
incoming inventory of 211 contains both a right arrival and a left arrival;
it does not supply a uniform immediate-left tag. This distinction is the
reason the stronger invariant helps.

\subsection{Four whole-tape identities}

Set the old head to coordinate zero by translation. The following table
lists all exceptional old-state/read cases. The stated cells are known;
every other cell is arbitrary and remains unread by the replacement block.

\begin{table}[ht]
\centering
\begin{tabular}{@{}cclcl@{}}
\toprule
Old state & Read & Needed context & New steps & Final state/head\\
\midrule
172 & 0 & \(t(-2)=1,\ t(-1)=0\) & 7 & \(67,-1\)\\
172 & 1 & none & 1 & \(184,+1\)\\
211 & 0 & \(t(-1)=1\) & 3 & \(157,+1\)\\
211 & 1 & none & 1 & \(184,+1\)\\
\bottomrule
\end{tabular}
\caption{Each endpoint is exactly the old one-step endpoint after applying
\(\rho\). Both machines preserve the complete tape in these cases.}
\label{tab:cases}
\end{table}

For old 211/read-0, Lemma~\ref{lem:211} gives the tape fragment 10 at
coordinates \((-1,0)\). The new path is
\[
211@0\longrightarrow217@{-1}\longrightarrow215@0\longrightarrow157@1.
\]
The successive rows are new 211/read-0, unchanged 217/read-1,
and unchanged 215/read-0. All writes preserve the scanned bit. This
is the old \(211/0=\act{0}{R}{157}\) action after three steps.

For old 172/read-0, the fragment at \((-2,-1,0)\) is 100. The complete
new trace is displayed explicitly; no old version of a changed helper
is used during this trace.

\begin{table}[ht]
\centering
\begin{tabular}{@{}rcrclcr@{}}
\toprule
Step & State & Head & Read & Executed action & Next state & Next head\\
\midrule
1 & 211 & 0 & 0 & \(\act{0}{L}{217}\) & 217 & $-1$\\
2 & 217 & $-1$ & 0 & \(\act{0}{R}{163}\) & 163 & 0\\
3 & 163 & 0 & 0 & \(\act{0}{L}{211}\) & 211 & $-1$\\
4 & 211 & $-1$ & 0 & \(\act{0}{L}{217}\) & 217 & $-2$\\
5 & 217 & $-2$ & 1 & \(\act{1}{R}{215}\) & 215 & $-1$\\
6 & 215 & $-1$ & 0 & \(\act{0}{R}{157}\) & 157 & 0\\
7 & 157 & 0 & 0 & \(\act{0}{L}{67}\) & 67 & $-1$\\
\bottomrule
\end{tabular}
\caption{The seven-step simulation of old 172/read-0. The tape is 100 on
the three displayed cells throughout; its arbitrary exterior is untouched.}
\label{tab:seven}
\end{table}

The old action \(172/0=\act{0}{L}{67}\) has exactly that endpoint.
The common read-one action \(\act{1}{R}{184}\) handles the remaining
two cases in a single step. None of the four blocks enters halt internally.

For every old row outside states 172 and 211, the new action writes the
same bit, moves in the same direction, and enters the \(\rho\)-image of
the old destination. These 414 ordinary rows therefore simulate in one step.

\subsection{Simulation, reflection of halting, and the clock}

\begin{theorem}[Uniform preservation of halting]
\label{thm:all-tape}
From state 0, the exact principal U209 and U208 tables halt on exactly
the same initial binary tapes. At matching points their tapes and head
positions agree. If the old computation halts after \(n\) steps, the new
one halts after
\[
\tau(n)=n+6A(n)+2B(n),
\]
where \(A(n)\) counts old 172/read-0 transitions and \(B(n)\) counts old
211/read-0 transitions among those \(n\) steps. In particular,
\(n\le\tau(n)\le7n\).
\end{theorem}
\begin{proof}
The initial configurations satisfy \(\mathcal R\), because \(\rho(0)=0\).
Induct on the number of old steps. In an ordinary state, the single new
step commutes with \(\rho\). In states 172 and 211, the old-machine
lemmas supply the relevant guards, and the appropriate block in
Table~\ref{tab:cases} returns to \(\mathcal R\). Summing the block
lengths gives the stated clock formula. All block lengths are positive.

If the old computation halts, the constructed finite new computation
reaches the corresponding halt. Conversely suppose the new machine
halts at time \(T\). If the old machine did not halt in its first \(T\)
steps, the induction would construct a new computation through
\(\tau(T)\ge T\) whose synchronized endpoints are working states and
whose internal blocks contain no halt. This contradicts halting at time
\(T\). Hence the old computation also halts.
\end{proof}

There is no circularity in the internal revisits to new 211 and new 163.
The lemmas concern only the old machine at old-step boundaries. The
seven-step new trace is executed directly, including the new 211/read-0
row on its second use. We do not assert the old invariant at each
transient configuration of the new machine.

For the PA version, code finite-support initial tapes and finite computation
histories by natural numbers. The predecessor facts and local traces are
fixed finite arithmetic verifications. Induction defines the synchronization
clock and proves the same endpoint relation; the finite \(T\)-step
reflection argument proves the reverse implication. Thus the blank instance
is an ordinary PA derivation in principle, without quantifying over arbitrary
sets of integers inside PA. The external all-tape theorem is stronger in
scope than the particular arithmetic blank-input statement required here.

\section{Discovery method and why the earlier screen missed it}
\label{sec:discovery}

The existing Joint research programs examine direct aliases and folds
supported by immediate incoming-edge tags. The published negative runs
are bounded screens of restricted templates, not proofs of minimality.
The two states used here agree on read-one but differ on read-zero:
\[
\begin{array}{c|cc}
 &0&1\\\hline
172&0L67&1R184\\
211&0R157&1R184.
\end{array}
\]
Immediate tags describe a zero to the right after a backward arrival at
211; they do not determine the bit needed by a new backward selector.
Looking a few transitions further back, while conditioning on reading
zero now, rules out two incoming branches and forces the missing one.

The included discovery program constructs finite backward symbolic
histories from the exact old table. It uses incoming write/read consistency
to discard impossible histories and retain forced bits. Candidate
replacement rows are screened on concrete examples to reject obvious
failures, then accepted only by exhaustive symbolic execution of every
remaining guarded case with arbitrary exterior. Random screening is used
only for rejection; it is not a correctness certificate.

The cleaned discovery script, with history depth 3 and macro bound 7,
rediscovers the two symmetric orientations of this fold. The final proof
was then simplified to Lemmas~\ref{lem:172} and~\ref{lem:211}; the release
checker does not trust the discovery engine, its random screen, or its
general backward-history implementation. It reads the old and new tables
and checks the short, fixed premises independently.

The invariant is necessary for this argument. For example, at an arbitrary
configuration of old 211 with local fragment 100 at coordinates
\((-2,-1,0)\), old 211/read-0 goes to \(157@1\); the new control follows
the seven-step behavior ending at \(67@{-1}\). Such a configuration is
excluded by Lemma~\ref{lem:211} on old runs from state 0. The result is
therefore a reachability-sensitive simulation, not equality of the two
state functions on unrestricted configurations.

\section{The logical calibration retained from U209}
\label{sec:calibration}

\subsection{The unary theory and its restriction}

The supplied unary manuscript~\cite{incoming} uses classical equality-free first-order
logic on nonempty domains, with one binary predicate \(E\). Write
\[
\Sub(x,a):=\forall z\,(E(z,x)\to E(z,a)).
\]
For each formula \(\varphi\) with free variables contained in \(\{x\}\),
the theory U contains
\[
\forall a\,\exists b\left[
\forall x\,(E(x,b)\to\varphi(x))\ \land\
\forall x\,(\Sub(x,a)\to(\varphi(x)\to E(x,b)))\right].
\tag{U}
\]
The source parameter \(a\) is not allowed free in \(\varphi\).
There is no bound on bound-variable count, formula size, or quantifier
complexity. No extensionality axiom is assumed. In particular, replacing
\(\varphi(x)\) by \(\Sub(x,a)\land\psi(x)\) would violate the stated
unary restriction, so ordinary unrestricted separation cannot simply be
inserted into the lower interpretation.

\subsection{What the manuscript proves}

The manuscript contains a substantial ordinary proof of the U/PA
consistency equivalence. Our review found no concrete defect in the
following delicate points; a new formal endpoint was not produced.

Coextensionality identifies objects with the same membership profile.
Numerical objects are saturated, transitive, linearly ordered in the
relevant membership sense, and satisfy least/greatest endpoint conditions
for represented intersections at initial cutoffs. These definitions do
not assume external finiteness or external well-foundedness.

The lower construction first proves an existing nonempty numerical object
has a greatest member that supplies its predecessor. It then proves
induction \emph{before} proving total successor. A unary cover selects
failures inside a numerical source. If there is an internal failure,
choose its least class; if there is none but the source itself fails,
the source is the least failure. In either case its predecessor contradicts
the induction step. This second case avoids circular reliance on an
already available successor.

Parameters are internalized through a legal unary predicate of the form
\[
Q(v):=\forall\vec p\,
\bigl(\mathrm{BaseStep}(\theta,\vec p)\to\theta(v,\vec p)\bigr).
\]
Least-failure arguments then recover exact bounded comprehension with
numerical parameters. A least counterexample tuple is bound inside the
unary matrix; an independent support contains the selected objects as
subsets. The proved invariance of the target formula in its numerical
parameters makes equivalent minimal tuples select the same profile.
This is a formula-specific syntactic construction, not a truth
predicate for all formulas.

Total successor and exact powers of numerical sources follow. Saturated singleton and pair
profiles provide relation coding modulo coextensionality. Finite graph
histories define addition and multiplication. Old graph witnesses are
existentially bound inside the construction; only numerical values remain
free parameters. Pointwise uniqueness and numerical bounds are established
before the relevant graph extensions. The translated induction schema
gives a full interpretation of PA.

In the other direction, PA interprets \(E(x,a)\) as the \(x\)-th binary
digit of \(a\). PA proves that \(\Sub(x,a)\) entails \(x\le a\).
For any fixed translated formula, induction constructs its characteristic
code on that bounded interval. This verifies even the unrestricted weak
covering schema. The same interpretation satisfies global irreflexivity.
These arguments are available in PA models, rather than only in the
standard natural numbers.

Uniform renaming, substitution, relativization, and finite derivation
templates provide the advertised proof translations. Their arithmetic
verification uses induction on formulas and finite derivations. The
manuscript does not deliver an implemented arbitrary-proof translator or
a checked Lean/Rocq/MM0 endpoint; that is an explicit formalization limit.

\subsection{What the literal machine actually enumerates}

The U209 lineage does not compile the unary ZIP's separate positive
dyadic stream. It searches an inherited native numerical calculus for
unrestricted weak covering \(W\) plus global irreflexivity \(\IR\).
Its target is native code 1, and code 0 is a truth alias. The positive
stream has a different target, code 0, and a different codec. A small
constructor count in that stream cannot be read as a state bound for
the literal U209 or U208 tables.

The consistency translations form the cycle
\[
\Con(\PA)\ \Longrightarrow\ \Con(W+\IR)\
\Longrightarrow\ \Con(U)\ \Longrightarrow\ \Con(\PA).
\]
The machine's additional native-calculus bridge is substantive. The
inherited numerical constructors use a Cantor pairing \(P\), implication
\(4P\), membership \(4P+1\), universal quantification \(4P+2\), and
auxiliary equality \(4P+3\). The alias \(I(0,0)=0\) and the target
\(R(0,0)=1\) require explicit treatment. Auxiliary equality is calibrated
through membership indiscernibility rather than importing unavailable
primitive equality rules.

The role-based covering emitter uses one integer both as matrix code and
output label. Its numerical bounds ensure freshness for positive codes;
the zero truth alias and source/output collision are handled separately.
Padding recovers every required cover instance. We inspected the retained U222 native
calibration and the later U214 role-emitter argument, rather than inferring
these properties from the final state count.

\subsection{Startup and proof-search completeness}

Blank startup enters an interior axiom-generation suffix. Its existing
cached value must not simply be treated as a theorem. The inherited
argument produces nine theorem-valued heads above an opaque stack tail,
then uses a diagonal credit bound to show the unsafe tail is never consumed.
A ghost stack length is retained because the numerical pairing has a zero
alias. Strong source compilation, including its raw-zero exceptions,
gives cofinally many successful refutation candidates and termination of
unsuccessful candidate evaluations. These are the source-to-native
soundness and completeness dependencies behind the blank-input endpoint.

The exact Joint release and its full nested ancestry are retained as
\art{baseline/U209_ConPA_Joint_2026-10-29.zip}. The unary-theory archive
is also retained separately. Some historical Rocq files in that ancestry
are reference excerpts with external imports; their presence is not a new
runnable U208 formal verification project.

The K215 ZFC formalizations and the K230 account of Theory K are useful
reference points~\cite{k215pa,k215formal,k230,m254}. They establish
different literal-table endpoints. The present seven-state numerical
difference from 215 does not by itself compare the intrinsic complexity
of PA and ZFC encodings.

\section{Verification of the delivered principal candidate}
\label{sec:verification}

\subsection{Independent local certificate}

\art{code/verify_fold.py} imports none of the discovery, source compiler,
or universal-runtime modules. It pins the exact parent, reads the actual
candidate, verifies the complete predecessor inventories, checks all
ordinary transition correspondences, and executes the four exceptional
blocks with an opaque exterior. A read outside the supplied finite
context fails; equality of the final finite cells therefore proves equality
of the full arbitrary tape.

Eight deliberately corrupted constructions are rejected, including wrong
first and second entry bits, a wrong deeper predecessor, a wrong helper
direction or write, and a wrong selector continuation. The local checker
passes under ordinary Python and under \texttt{-O}; its required conditions
use explicit exceptions rather than removable assertions. The discovery
script is not invoked by the release verifier.

\subsection{Unchanged source interfaces and startup}

The source bytes are identical to Joint's. All 3,908 internal source
prefixes were traversed against both tables, with endpoints commuting
with \(\rho\). The maximum prefix traversal length changes from 49 to
50 literal steps; no source address changes.

The complete 48,140-step blank startup was executed independently on
both tables. The entire transition trace, final head, state, and finite
support of the tape agree. No edited or deleted row is visited during
that prefix. The endpoint is native state 84, head \(-2\), source
ordinal 136, address \texttt{001100110101010}, and logical registers
\[
[1,0,13,0,0,1,0,0,1],
\]
with 43 spare registers all equal to 1. The full canonical frame has
120 one-cells and zero exterior, recorded in the startup receipt.
This finite check verifies initialization; it is not evidence of
unbounded nonhalting by itself.

\subsection{Fresh universal primitive replay}

The inherited native-label induction engine was run on the final 208-state
table. For each source primitive it checks the complete tape endpoint
with arbitrary values of nine primary registers, an arbitrary natural
gap, and an arbitrary nonempty finite spare-register list. Conditional
decrement has separate zero and positive cases, giving
\(3909+1175=5084\) contracts.

The engine's inference rules are literal transition steps, induction over
a unary scan length for a corresponding self-loop, and structural induction
over a finite list after checking its one-block transformation. Four
one-block schemas are rederived from the literal table. Durations are
affine expressions whose coefficients are checked nonnegative and whose
constant terms are positive. The minimum constant is 5.

\begin{table}[ht]
\centering
\begin{tabular}{@{}lr@{}}
\toprule
Principal U208 check & Result\\
\midrule
Working states / complete transition rows & 208 / 416\\
Graph-reachable working states & 208\\
Exact exceptional local cases & 4\\
Other old rows commuting with projection & 414\\
Rejected local corruptions & 8\\
Source primitive sites / ordinal control arrows & 3,909 / 5,083\\
Internal source-prefix interfaces & 3,908\\
Universal primitive contracts & 5,084\\
Symbolic inference applications & 1,251,865\\
Literal transition rows covered & 416 / 416\\
Reproved finite-list block schemas & 4\\
Exact common startup steps & 48,140\\
\bottomrule
\end{tabular}
\caption{Verification facts for the principal candidate. Symbolic contracts
are unbounded mathematical identities, not 5,084 sampled register tuples.}
\label{tab:verification}
\end{table}

The Python induction engine is inherited, readable, and explicitly part
of the trusted executable evidence. It is not a proof-assistant kernel.
The local proof in Section~\ref{sec:proof} transfers the parent's theorem
independently of rerunning this engine; the direct replay is valuable
because it checks the delivered table against the unchanged source as well.

\section{A second 208-state table from Compatibility}
\label{sec:alternative}

The same new idea applies to the distinct Compatibility U209 table.
Delete 172, redirect 116/read-0 to 211, and use
\[
\delta(211,0)=\act{0}{L}{214}.
\]
The relevant role correspondence is
\begin{center}
\begin{tabular}{@{}lrrr@{}}
\toprule
 & Folded predecessor & New backward helper & Return selector\\
\midrule
Joint & 163 & 217 & 215\\
Compatibility & 216 & 214 & 163\\
\bottomrule
\end{tabular}
\end{center}
In Compatibility, 172 is still entered through 116/read-0, now with
116 entered through 163/read-1. State 211 is entered from 216; a backward
arrival reading zero forces the 214/read-0 predecessor, whose incoming
217/read-1 supplies the one. The same seven-, three-, and one-step
identities follow with the listed labels.

\art{code/check_alternative.py} verifies this as an explicit label
instantiation of the local theorem. This label correspondence is not
a claim that the two complete parent tables are isomorphic. The alternative
dense table is \art{machine/u208_compatibility.tm}, SHA-256
\begin{center}\small
\path{9865dcd7bd44deece098c16dcec7b90a96cf111592cbd37e3531b50f7b09276c}.
\end{center}
Its direct local preservation theorem is checked; a separate fresh full
primitive replay for this supplementary table was not performed. The
principal table remains the one with the full new runtime and startup receipts.

\section{Limits, excluded shortcuts, and further work}
\label{sec:limits}

The specific advance over the supplied proposals is the conditional
three-transition entry argument and the use of an already present decoder
to distinguish two local contexts. The final modification changes only
two retained rows and removes one state. No new logical axiom system,
proof-search scheduler, source reassociation, or register representation
is required for this saving. Novelty is claimed relative to the supplied
archives; no broader historical priority claim is made.

A larger search followed the same family of templates to history depth 8
and replacement-block bound 31. It found the same two orientations.
A follow-up search on U208 examined 448 oriented row-compatible pairs
and 369,152 replacement-row choices without another success; a direct
alias search was also unsuccessful. These are restricted experimental
negative results. They establish neither a lower bound of 208 nor
exhaustion of possible representation changes.

Two tempting theory-level shortcuts are already ruled out by the input
research. Deleting either freshness gate alone from the unchanged
covering emitter has small-code counterexamples. Deleting the whole
``dominated'' propositional row is also invalid as a route to the same
calibration: Plain constructs a Boolean valuation with implication
interpreted as equality, membership false, auxiliary equality true,
and quantifiers acting identically on truth values. All remaining
axioms are designated and modus ponens/generalization preserve that
property, while the target is not designated. PA can verify underivability
of the target in the reduced calculus. If PA is consistent, that calculus
cannot retain a PA-provable equivalence to \(\Con(\PA)\), by second
incompleteness. These are inherited obstructions, not new discoveries here.

The most direct next verification task is to encode the two entry lemmas
and the generic positive-block simulation theorem in a small formal
framework, then instantiate them with the sealed table. That would
formalize the new reduction alone. A fully checked Con(PA) endpoint
still requires formalization of the inherited U/native/source calibration.
For further state reduction, guarded history information is a useful
additional search dimension; the present bounded searches do not decide
what can be achieved with longer contexts or more global transformations.

\section{Archive guide and reproduction}
\label{sec:reproduce}

The principal deliverables are this PDF and its complete LaTeX source,
the dense and native U208 tables, and deterministic offline verification.
\begin{center}\small
\begin{tabularx}{\linewidth}{@{}lX@{}}
\toprule
Path & Contents\\
\midrule
\art{report.tex}, \art{report.pdf}, \art{paper/} & Article, generated full table, and provenance appendix.\\
\art{machine/} & Principal and alternative literal tables; unchanged source; dense/native map.\\
\art{code/} & Builder, independent local checker, standalone comparison, source/startup audit, and universal runtime wrapper.\\
\art{evidence/} & Actual local, comparison, startup, and universal-contract receipts.\\
\art{comparison/} & Exact tables/source maps, notices, and reports from the six incoming packages.\\
\art{baseline/} & Complete sealed Joint parent ZIP, unary theory ZIP, exact parent tables, and source manifest.\\
\art{research/} & Clean discovery script, discovery receipt, and detailed comparison/theory audit notes.\\
\art{licenses/}, \art{NOTICE.txt} & Preserved attribution and inherited terms.\\
\art{SHA256SUMS} & Hashes of the other sealed release files.\\
\bottomrule
\end{tabularx}
\end{center}

Use Python 3.10 or later, with its standard library and assertions enabled:
\begin{verbatim}
python3 -B reproduce.py --output /tmp/u208-verification.json
\end{verbatim}
The script verifies the manifest, reconstructs the exact machine bytes
in a temporary directory, checks the independent local proof both normally
and under \texttt{-O}, checks the alternative, compares all six snapshots,
replays startup, and freshly rederives all 5,084 contracts. The sealed
release is not modified. The runtime phase takes a few minutes on the
preparation environment; elapsed time is not part of the certificate.

\begin{verbatim}
python3 -B reproduce.py --quick
\end{verbatim}
explicitly skips the universal runtime phase and reports that omission.
The optional \texttt{--parent-replay} also unpacks and replays the complete
retained Joint release. If Poppler's \texttt{pdftotext} is present, comparison
freshly extracts the PDF-only table. Otherwise it verifies the delivered
extracted table and explicitly records that extraction was not rerun;
the standalone comparison's \texttt{--require-pdf-extraction} makes that
tool mandatory. Preparation included fresh extraction.

To rerun discovery separately:
\begin{verbatim}
python3 research/discover_folds.py baseline/u209.native.tm \
  --output /tmp/u208-discovery.json
\end{verbatim}
Discovery is optional and is not a prerequisite for verifying the fixed
construction. To rebuild the article in a working copy, run
\texttt{pdflatex -interaction=nonstopmode -halt-on-error report.tex} twice.
TeX versions and PDF metadata may change PDF bytes; the machine bytes are
fixed and checked independently of typesetting.

\appendix
\clearpage
\section{Complete dense transition table}
\label{app:table}

This is the principal \art{machine/u208.tm}. Each entry \(wDz\) means
write \(w\), move \(D\), and enter \(z\). Start in 0 on the all-zero
tape; 208 is the separate halt label. Two state groups are printed side
by side to save space. Every row is generated directly from the sealed
literal table, not transcribed manually.
\begingroup\footnotesize
\setlength{\tabcolsep}{12pt}
\begin{longtable}{@{}rll@{\hspace{24pt}}rll@{}}
\toprule
State & Read 0 & Read 1 & State & Read 0 & Read 1\\
\midrule\endfirsthead
\toprule
State & Read 0 & Read 1 & State & Read 0 & Read 1\\
\midrule\endhead
\midrule\multicolumn{6}{r}{\small Continued on the next page}\\\endfoot
\bottomrule\endlastfoot
0 & 1R92 & 1R0 & 104 & 0R21 & 1R120\\
1 & 0L1 & 1L2 & 105 & 0R21 & 0L10\\
2 & 0L3 & 1L2 & 106 & 0R61 & 1R42\\
3 & 0L5 & 1L2 & 107 & 0R121 & 1R8\\
4 & 1R6 & 0L4 & 108 & 0R101 & 1R122\\
5 & 0L9 & 1L2 & 109 & 0L207 & 0L95\\
6 & 0R6 & 0R10 & 110 & 0R124 & 1R104\\
7 & 0R189 & 1R159 & 111 & 0R115 & 1L72\\
8 & 0R14 & 1R13 & 112 & 0R194 & 1R180\\
9 & 0L15 & 1L2 & 113 & 0R160 & 1R183\\
10 & 0L16 & 0R10 & 114 & 0R148 & 1R146\\
11 & 0L26 & 1R128 & 115 & 0R11 & 1R129\\
12 & 0R19 & 1R20 & 116 & 0L131 & 0L1\\
13 & 0L89 & 1L31 & 117 & 0R39 & 1R72\\
14 & 0L11 & 0R22 & 118 & 0R61 & 1R132\\
15 & 0L23 & 1L2 & 119 & 0R61 & 1R133\\
16 & 1L1 & 0L1 & 120 & 0R70 & 0L16\\
17 & 0R24 & 0R32 & 121 & 0R14 & 1R132\\
18 & 0R98 & 1R113 & 122 & 0R21 & 1R134\\
19 & 0R28 & 1R29 & 123 & 0R151 & 1R188\\
20 & 0R30 & 1R31 & 124 & 0R137 & 1L20\\
21 & 0R24 & 0R59 & 125 & 0R130 & 1R179\\
22 & 0R33 & 1R22 & 126 & 0R203 & 1R156\\
23 & 0L34 & 1L2 & 127 & 0R75 & 1R175\\
24 & 0L35 & 0R0 & 128 & 0R75 & 1R205\\
25 & 0R36 & 1R25 & 129 & 0R202 & 1R157\\
26 & 0R37 & 1R38 & 130 & 0R12 & 1R141\\
27 & 0R39 & 0L40 & 131 & 0R30 & 1R69\\
28 & 0R41 & 1R42 & 132 & 0R39 & 1L124\\
29 & 0R43 & 1R44 & 133 & 0R153 & 0L16\\
30 & 0R45 & 1R46 & 134 & 0R154 & 0L16\\
31 & 0R21 & 1R47 & 135 & 0R88 & 1R136\\
32 & 0R48 & 1R32 & 136 & 0R196 & 1R140\\
33 & 0R32 & 1R33 & 137 & 0R24 & 0R48\\
34 & 0L74 & 1L2 & 138 & 0R99 & 1R37\\
35 & 1R35 & 1L1 & 139 & 0R171 & 1R148\\
36 & 0R50 & 1R36 & 140 & 0R96 & 1R159\\
37 & 0R17 & 1R111 & 141 & 0R174 & 1R165\\
38 & 0R13 & 1R52 & 142 & 0R12 & 1R159\\
39 & 0L11 & 0R36 & 143 & 0R88 & 1L165\\
40 & 0L4 & 1L53 & 144 & 0R88 & 1R158\\
41 & 0R24 & 0R25 & 145 & 0L152 & 1L163\\
42 & 0R13 & 0L16 & 146 & 0R84 & 1R88\\
43 & 0R55 & 1R56 & 147 & 0R159 & 1R12\\
44 & 0R57 & 1R58 & 148 & 0R83 & 1R143\\
45 & 0R24 & 0R54 & 149 & 0R144 & 1R155\\
46 & 0R60 & 1R61 & 150 & 0L66 & 1L141\\
47 & 0R62 & 0L16 & 151 & 0R150 & 1R169\\
48 & 0R25 & 1R48 & 152 & 0R159 & 1R114\\
49 & 0L35 & 0R63 & 153 & 0R97 & 1R97\\
50 & 0R16 & 1R50 & 154 & 0R49 & 0R32\\
51 & 0R147 & 1R143 & 155 & 0R191 & 1R187\\
52 & 0R67 & 1R68 & 156 & 0L194 & 1R194\\
53 & 0L40 & 1L53 & 157 & 0L186 & 1R96\\
54 & 0R22 & 1R54 & 158 & 0R167 & 1R178\\
55 & 0L11 & 0R50 & 159 & 0R19 & 1R170\\
56 & 0R60 & 0L16 & 160 & 0R164 & 1R186\\
57 & 0R69 & 1R70 & 161 & 0R159 & 1R96\\
58 & 0R61 & 1R71 & 162 & 0R125 & 1R164\\
59 & 0R54 & 1R59 & 163 & 0R97 & 1R162\\
60 & 0R49 & 0R25 & 164 & 0L189 & 1R142\\
61 & 0L28 & 1L72 & 165 & 0R146 & 1R171\\
62 & 0L117 & 1L30 & 166 & 0R96 & 1R123\\
63 & 0R73 & 1R63 & 167 & 0R161 & 1R143\\
64 & 0R97 & 1R150 & 168 & 0R148 & 1R87\\
65 & 1R184 & 1L53 & 169 & 0R161 & 1R96\\
66 & 0R144 & 1R145 & 170 & 0R180 & 1R122\\
67 & 0R76 & 1R208 & 171 & 0L192 & 1R60\\
68 & 0L196 & 1R31 & 172 & 0R145 & 1R88\\
69 & 0R41 & 1R62 & 173 & 0R152 & 1R201\\
70 & 0L89 & 1R78 & 174 & 0R140 & 1R84\\
71 & 0R55 & 1L79 & 175 & 0R174 & 1R75\\
72 & 0R80 & 1R81 & 176 & 0R126 & 1R168\\
73 & 1R0 & 1R63 & 177 & 0R138 & 1R116\\
74 & 0R74 & 1R82 & 178 & 0R97 & 1R51\\
75 & 0L145 & 1L141 & 179 & 0R167 & 1R51\\
76 & 0R85 & 1R86 & 180 & 0R17 & 1L131\\
77 & 0R20 & 1R38 & 181 & 0R143 & 1R149\\
78 & 0R89 & 1L132 & 182 & 0R76 & 1R116\\
79 & 0L108 & 1R56 & 183 & 0R106 & 1R124\\
80 & 0R49 & 0R33 & 184 & 1L40 & 1L109\\
81 & 0R91 & 0L42 & 185 & 0R174 & 1R156\\
82 & 0R177 & 1R77 & 186 & 0R181 & 1R182\\
83 & 0R96 & 1R12 & 187 & 0R66 & 1R204\\
84 & 0R88 & 1R97 & 188 & 0R140 & 1R197\\
85 & 0R206 & 1R21 & 189 & 0R84 & 1R190\\
86 & 0R116 & 1R100 & 190 & 0R157 & 1R185\\
87 & 0R173 & 1R166 & 191 & 0R64 & 1R7\\
88 & 0R102 & 1R103 & 192 & 0R150 & 1R127\\
89 & 0R80 & 1R79 & 193 & 0R18 & 1R94\\
90 & 0R101 & 1R195 & 194 & 0L200 & 1R172\\
91 & 0R105 & 1R106 & 195 & 0R104 & 1R96\\
92 & 0R65 & 0R0 & 196 & 0R90 & 1R88\\
93 & 0R27 & 1R107 & 197 & 0R139 & 1R199\\
94 & 0R182 & 0L207 & 198 & 0R150 & 1R112\\
95 & 1L109 & 1L95 & 199 & 0R201 & 1R200\\
96 & 0R19 & 1R110 & 200 & 0R156 & 1R198\\
97 & 0R60 & 1R60 & 201 & 0L123 & 1R167\\
98 & 0R68 & 1R135 & 202 & 0R196 & 1R28\\
99 & 0R93 & 1R108 & 203 & 0R163 & 1R192\\
100 & 0R105 & 0L207 & 204 & 0R146 & 1R176\\
101 & 0R117 & 1R118 & 205 & 0R161 & 1R193\\
102 & 0R71 & 1R58 & 206 & 0R45 & 1R100\\
103 & 0R57 & 1R119 & 207 & 0L4 & 0L207\\
\end{longtable}
\endgroup


\clearpage
\section{Source identities and evidence provenance}
\label{app:provenance}
\subsection*{Principal machine and unchanged source}

\textbf{u208.tm}\par

\begin{small}\path{a2fee9e41311c0d43aaac871ea411be24eada17e6b3633c55f40b437101144e1}\end{small}\par

\textbf{u208.native.tm}\par

\begin{small}\path{b348a466a2c316511e0037ea3824de29d4c6ad37338404fbc98a3968fdb25135}\end{small}\par

\textbf{sites.json}\par

\begin{small}\path{cddbed7f20e7a8b085b7ca0e76811d07b021522375643b5d1134eec0d2d23164}\end{small}\par

\subsection*{Retrieved incoming archives}

The hashes below identify the complete original ZIP bytes, including nested artifacts.

Original Git blob hashes, sizes, and pinned download URLs accompany them in the JSON manifest.

\noindent\textbf{\path{U209_ConPA_2026-10-29.zip}}\par

\small\path{7fd484091428643f211ea55d8750c127734705c0f91cd549b74fdc2c0d598f9f}\normalsize\par

\noindent\textbf{\path{U209_ConPA_Compatibility_2026-10-29.zip}}\par

\small\path{35c9c1c9f5526f7f85200fd2b2f872ade9c05655b6722d40d9dc5232958a8ec4}\normalsize\par

\noindent\textbf{\path{U209_ConPA_Compatible_Folds_2026-10-29.zip}}\par

\small\path{afee49ed265309f84e69eb0d92a5cfba929377381e570a2c562fa2f5702d6f4d}\normalsize\par

\noindent\textbf{\path{U209_ConPA_Composed_2026-10-29.zip}}\par

\small\path{5ad6da480e4fcf22ad020659070f487d73efbd8aefe9a20fcf5f00195853b9d0}\normalsize\par

\noindent\textbf{\path{U209_ConPA_Joint_2026-10-29 (1).zip}}\par

\small\path{981218c643c6cc3a62afaa93e9c6cf0f73dce203e1b344d33ee59e9a9a6bd91e}\normalsize\par

\noindent\textbf{\path{U209_ConPA_Joint_2026-10-29.zip}}\par

\small\path{9dd2199baffa0244da981384c81920afdccc353509c3c8745b28e95a93afc7a1}\normalsize\par

\noindent\textbf{\path{pa_unary_power_cover.zip}}\par

\small\path{0fcb1117a5353766a3366bab04e56b4efbd0607833cdc73ffb34b67e6808b0ff}\normalsize\par

\subsection*{Retained executable evidence}

The new local certificate is \path{evidence/independent_fold.json}; the optimized replay is \path{evidence/independent_fold_optimized.json}. The exact source/startup receipt is \path{evidence/source_startup.json}. Fresh initial universal data are summarized in \path{evidence/u208_runtime_initial.json}, with all case records in \path{evidence/u208_runtime_cases.json.gz}. The standalone six-way comparison is \path{evidence/six_way_comparison.json}; its PDF extraction was freshly executed. Original-parent full replays are recorded in \path{evidence/joint_replay.json} and \path{evidence/compatibility_replay.json}.

The new runtime wrapper rederives its case records and requires equality with the sealed records. These records are supporting executable evidence; the fixed local proof does not assume the search algorithm is sound. The archive manifest binds the code and machine bytes supplied for review.


\begin{thebibliography}{9}
\bibitem{incoming}
Vladimir Reshetnikov, ProveIt, six U209 archives and
\texttt{pa\_unary\_power\_cover.zip}, snapshot
\texttt{343da82d7384d10ec9fd1c17361fc923b2d1d117}, 29 October 2026.
\url{https://github.com/VladimirReshetnikov/ProveIt/tree/343da82d7384d10ec9fd1c17361fc923b2d1d117/docs/incoming}.
Exact filenames and archive hashes are listed above and in the source manifest.
\bibitem{m254}
Jeff's busy-beaver repository, \texttt{ZF/m254}, branch \texttt{zf-254}.
\url{https://github.com/jeff87654/busy-beaver/tree/zf-254/ZF/m254}.
\bibitem{k215formal}
ProveIt K215 verification: Rocq, Lean, and audit materials.
\url{https://github.com/VladimirReshetnikov/ProveIt/tree/343da82d7384d10ec9fd1c17361fc923b2d1d117/Computability/BusyBeaver/Research/ZFC/k225-formal/verification/k215}.
\bibitem{k215pa}
Elliot Glazer, \texttt{k215-pa}, the literal K215/Con(ZFC) theorem and
MM0 verification receipts.
\url{https://github.com/elliotglazer/k215-pa}.
\bibitem{k230}
Elliot Glazer, \texttt{k230-rocq}, README, ``The theory K.''
\url{https://github.com/elliotglazer/k230-rocq/blob/main/README.md#the-theory-k}.
\end{thebibliography}
\end{document}
